\subsubsection{Security model and a rotation-moment bound}
\label{sec:qef_moment_security}
{\color{red}
The security quantity considered here is the number of independently
encrypted quantum preparations needed to distinguish two specified images.
This is a sample-complexity statement, rather than a gate-complexity estimate
for exhaustive key search. The proof uses the original three QEF layers.
It does not assume that an adversary reconstructs an arbitrary unitary or
recovers every component of the original key.

\paragraph{Keys, ciphertexts, and the attack experiment.}
Write $A$ for the substitution-angle material produced by the nD-QHCS and
$B=(K^{\rm perm},K^{\rm ent})$ for the material produced by the
HG$\Phi$-QRNG. These are two source classes. The two components of $B$
are described by their joint distribution; independence is not inferred
from their output lengths. Each encryption uses a fresh secret key record,
and the receiver obtains the same record over an assumed secure channel.
Independent refresh between preparations is an explicit security
assumption, stronger than merely obtaining different numerical key values.
The experiment gives the adversary no additional side information
correlated with the target secret record; any claimed extension to
source leakage must bound that leakage explicitly.

Let $N=HW$, let $r$ denote the feature-register size, and let
$D=N2^r$ denote the dimension of the valid image subspace. Unused padded
position addresses are excluded from the ciphertext experiment. The image
states have the form
\begin{equation}
\label{eq:qef_security_image}
 |I\rangle=N^{-1/2}\sum_{k=0}^{N-1}|k\rangle|\chi_k\rangle,
 \qquad \langle\chi_k|\chi_k\rangle=1.
\end{equation}
The adversary is given the classical descriptions of two such images of
the same public dimensions. A uniform hidden bit $b$ selects $I_b$.
The adversary receives $q$ encryptions of that selected image, with
independently refreshed keys, and may perform any joint quantum
measurement. The publicly transmitted ciphertext is the quantum state,
not a classical list of all its amplitudes or a rendered image file.
Define
\begin{equation}
\label{eq:qef_averaged_cipher}
 \sigma_b=\mathbb E_{A,B}
   [U_{A,B}|I_b\rangle\langle I_b|U_{A,B}^{\dagger}].
\end{equation}
The joint ciphertext is $\sigma_b^{\otimes q}$. Its optimal binary
discrimination probability is \cite{Watrous2018QEF}
\begin{equation}
\label{eq:qef_operational_guess}
 P_{\rm guess}^{(q)}
 =\frac12+\frac14\left\|
       \sigma_0^{\otimes q}-\sigma_1^{\otimes q}\right\|_1.
\end{equation}
This experiment does not give repeated access to one fixed-key unitary.
It also does not cover key reuse, arbitrary reference-entangled chosen
inputs, authentication, or publication of simulator amplitude records.

\paragraph{Use the actual circuit structure.}
For the intended pixel-dependent substitution, write
\begin{equation}
 R_A=\sum_k |k\rangle\langle k|\otimes R_{A,k},\qquad
 R_{A,k}=\bigotimes_{j=1}^r T_{k,j},\quad
 T_{k,j}=R_z(\theta_{k,j}^z)R_y(\theta_{k,j}^y)R_x(\theta_{k,j}^x).
\end{equation}
Position-controlled CNOT/CZ gates leave each position basis label
unchanged. Thus, for suitable feature unitaries $V_B(t)$,
\begin{equation}
 E_B=\sum_t |t\rangle\langle t|\otimes V_B(t),\qquad
 U_{A,B}=E_BP_BR_A
 =\sum_k |\pi_B(k)\rangle\langle k|
       \otimes V_B(\pi_B(k))R_{A,k}.
\end{equation}
The permutation and entanglement keys may have any joint distribution
in the theorem below. In particular, a uniform distribution over all
$N!$ permutations is not needed for the image experiment in
Eq.~\eqref{eq:qef_security_image}. The permutation must nevertheless be
implemented reversibly with its work registers correctly cleaned.

\paragraph{An explicit angle reference law.}
The following law is an analytical reference, not an experimentally
certified property of the nD-QHCS. For an integer $L\geq2$, let each
Euler angle be uniform on
\begin{equation}
 \Theta_L=\{2\pi s/L:s=0,\ldots,L-1\}.
\end{equation}
Conditioned on $B$, the $3r$ angles within a pixel are independent,
and the angle tuples belonging to any two different pixels are
independent. Pairwise independence of pixel tuples is sufficient;
independence of all $N$ tuples jointly is not required. The continuous
reference is the corresponding independent uniform law on $[0,2\pi)$.
These conditions must hold after conditioning on any public setup
information or side information used in the attack experiment.
The reference is a key-law assumption, not a consequence of deterministic
seed expansion. For $N\geq2$, two independent uniform angle tuples alone
require at least $6r\log_2L$ bits of conditional source entropy. The
continuous law is an ideal limiting model; a finite implementation must
use the finite-law expression and provide sufficient actual key material.

Set
\begin{equation}
\label{eq:qef_beta_definition}
 a_L=\frac{1}{L\sin(\pi/(2L))},\qquad
 \beta_L=a_L^{6r},\qquad
 \beta_{\infty}=\left(\frac{2}{\pi}\right)^{6r}.
\end{equation}
Uniform Euler coordinates are not Haar-distributed rotations. The
nonzero first moment is retained in the analysis below.

\begin{theorem}[Conditional image-discrimination sample bound]
\label{thm:qef_moment_sample_bound}
Under the stated angle reference law and independent key refresh, let
$N\geq2$, $c_N=1-1/N$, and
\begin{equation}
\label{eq:qef_information_budget}
 C_{N,r,L}=\ln\!\left(1+D\beta_L^2c_N\right).
\end{equation}
For every image state in Eq.~\eqref{eq:qef_security_image}, its average
ciphertext $\sigma$ satisfies
\begin{align}
 \frac12\|\sigma-I_D/D\|_1&\leq\beta_Lc_N,
 \label{eq:qef_one_image_distance}\\
 \mathsf D(\sigma\|I_D/D)&\leq C_{N,r,L},
 \label{eq:qef_cipher_relative_entropy}
\end{align}
where $\mathsf D$ uses natural logarithms. For any two such images,
achieving $P_{\rm guess}^{(q)}\geq p>1/2$ requires
\begin{equation}
\boxed{\displaystyle
 q\ \geq\
 \left\lceil
 \frac{\ln2-h(p)}{\ln(1+N2^r\beta_L^2(1-1/N))}
 \right\rceil,\qquad
 h(p)=-p\ln p-(1-p)\ln(1-p).}
\label{eq:qef_final_sample_bound}
\end{equation}
The continuous reference uses $\beta_\infty$ in the same expressions.
For $N=1$, the feature twirl makes the reference ciphertext exactly
maximally mixed, and no finite number of fresh encryptions distinguishes
the candidate images.
\end{theorem}

\begin{proof}
For one rotation axis $a$, averaging its conjugation over $\Theta_L$
dephases the Bloch components transverse to that axis. Composing the
independent $x$, $y$, and $z$ averages therefore gives
\begin{equation}
 \mathbb E[T X T^{\dagger}]=\frac{I_2}{2}\operatorname{Tr}X.
\end{equation}
For the first moment, a finite trigonometric sum gives
\begin{equation}
 \mathbb E[R_a(\theta)]
 =c_LI_2-is_L\sigma_a,
 \quad c_L=1/L,\quad s_L=L^{-1}\cot(\pi/(2L)).
\end{equation}
Each matrix on the right is $a_L$ times a unitary, since
$c_L^2+s_L^2=a_L^2$. Hence
$M=\mathbb E[R_{A,k}]=\beta_L^{1/2}W$ for a feature unitary $W$,
and the feature conjugation average is the completely depolarizing
channel. In the continuous limit,
\begin{equation}
 \mathbb E[T]=(2/\pi)^3 I_2,
 \qquad M=(2/\pi)^{3r}I_{2^r}.
\end{equation}
For example, the first identity follows by multiplying
$\mathbb E[R_a]=-i(2/\pi)\sigma_a$ in the order $z,y,x$.

Write $\rho=|I\rangle\langle I|=\sum_{k,l}|k\rangle\langle l|\otimes X_{kl}$,
and let $\Delta$ remove its off-diagonal position blocks. Conditional
pairwise independence gives
\begin{equation}
 \mathbb E[R_{A,k}X_{kl}R_{A,l}^{\dagger}]
 =\begin{cases}
 I_{2^r}\operatorname{Tr}(X_{kk})/2^r,&k=l,\\
 MX_{kl}M^{\dagger},&k\ne l.
 \end{cases}
\end{equation}
For Eq.~\eqref{eq:qef_security_image},
$\operatorname{Tr}X_{kk}=1/N$. Consequently the diagonal contribution
is $\tau=I_D/D$, which is invariant under every $E_BP_B$. The remaining
term is
\begin{equation}
 \sigma-\tau=\beta_L\,\mathbb E_B
  [F_B(\rho-\Delta\rho)F_B^{\dagger}],
 \quad F_B=E_BP_B(I\otimes W).
\label{eq:qef_off_diagonal_remainder}
\end{equation}
Every $F_B$ is unitary, so both trace norm and Hilbert--Schmidt norm
are bounded by the corresponding norms before this average.

The pinching map is an average over $N$ diagonal position-phase
unitaries, one of which is the identity. The triangle inequality gives
$\|\rho-\Delta\rho\|_1\leq2(1-1/N)$. Combining this with
Eq.~\eqref{eq:qef_off_diagonal_remainder} proves
Eq.~\eqref{eq:qef_one_image_distance}.
Moreover, because $\rho$ is pure and its position probabilities are
uniform,
\begin{equation}
 \|\rho-\Delta\rho\|_2^2
 =\operatorname{Tr}\rho^2-\operatorname{Tr}(\Delta\rho)^2
 =1-1/N.
\end{equation}
Thus $\|\sigma-\tau\|_2^2\leq\beta_L^2c_N$. Jensen's inequality applied to
the eigenvalues of $\sigma$ yields
\begin{align}
 \mathsf D(\sigma\|\tau)
 &=\ln D+\operatorname{Tr}(\sigma\ln\sigma)\\
 &\leq\ln(D\operatorname{Tr}\sigma^2)
 =\ln(1+D\|\sigma-\tau\|_2^2)
 \leq C_{N,r,L}.
\end{align}
For a uniform hidden bit and $q$ independent encryptions, the Holevo
information obeys
\begin{equation}
 \chi_q\leq\frac12\sum_{b=0}^1
 \mathsf D(\sigma_b^{\otimes q}\|\tau^{\otimes q})
 \leq qC_{N,r,L}.
\end{equation}
Indeed, subtracting $\chi_q$ from the first expression gives the
nonnegative relative entropy of the average ciphertext to
$\tau^{\otimes q}$. Holevo's theorem and binary Fano's inequality
\cite{Watrous2018QEF} imply that any guess with success at least $p$
requires mutual information at least $\ln2-h(p)$. Therefore
$\ln2-h(p)\leq qC_{N,r,L}$, proving the result. The same argument
allows any joint measurement and unlimited computational postprocessing.
\end{proof}

\paragraph{Numerical meaning and source certification.}
For the QIRBP feature count $r=11$, the continuous reference gives
$\beta_\infty\simeq1.13785812768644\times10^{-13}$.
Table~\ref{tab:qef_quantum_attack_resistance} evaluates
Eq.~\eqref{eq:qef_final_sample_bound} at $p=2/3$.
The numbers are necessary ciphertext-preparation counts under the
specified law; they are neither measured attack times nor bit-security
exponents. The logarithm of a sample count must not be labelled as
elementary-gate security.

\begin{table}[htbp]
\centering
\color{red}
\caption{{\color{red}Conditional QIRBP sample lower bounds for binary image
discrimination. Independent uniform Euler angles on $[0,2\pi)$,
fresh independent keys, $r=11$, and success $p=2/3$ are assumed.
The actual nD-QHCS/HG$\Phi$ source conditions remain uncertified.}}
\label{tab:qef_quantum_attack_resistance}
\small
\begin{tabular}{lrrr}
\toprule
Image & $N$ & $D=N2^{11}$ & Necessary preparations $q$\\
\midrule
Ashurbanipal & 65,536 & 134,217,728 & $\geq3.25904\times10^{16}$\\
Hammurabi & 50,394 & 103,206,912 & $\geq4.23831\times10^{16}$\\
Lamassu & 65,536 & 134,217,728 & $\geq3.25904\times10^{16}$\\
Babylon Lion & 786,432 & 1,610,612,736 & $\geq2.71583\times10^{15}$\\
\bottomrule
\end{tabular}
\end{table}

To relate the reference calculation to an actual generator, suppose
the resulting average ciphertext is within trace distance $\kappa$
of the reference ciphertext, uniformly over the allowed input images.
For example, a bound $\kappa$ on the total variation distance of the
actual joint finite angle/control-key law from its finite reference
law suffices. Comparing a finite law directly to a continuous law in
total variation is inappropriate; the reference must use the same
finite alphabet. Then the preceding proof gives
\begin{equation}
\label{eq:qef_source_robust_bound}
 C_{\kappa}=\ln\!\left(1+D
   [\beta_L\sqrt{1-1/N}+2\kappa]^2\right),\qquad
 q\geq\left\lceil\frac{\ln2-h(p)}{C_\kappa}\right\rceil.
\end{equation}
The trace-distance bound becomes
$\frac12\|\sigma-\tau\|_1\leq\beta_L(1-1/N)+\kappa$.
This follows from $\|\sigma-\sigma_{\rm ref}\|_2
\leq\|\sigma-\sigma_{\rm ref}\|_1\leq2\kappa$.
The source bound must include correlations between pixels, axes,
source parameters, both components of $B$, and any relevant side
information. Statistical image metrics or balanced one-bit marginals
do not certify this joint bound.

The mapping $\theta=\gamma\bmod(q,\pi)$ does not by itself establish
the reference law. Even if $q\bmod\pi$ were uniform, its angle range
would be $[0,\gamma\pi)$; the reference continuous range requires
$\gamma=2$, as well as the stated joint moment conditions. No measured
value of $\kappa$ or proof of those conditions for the original
generators is supplied here. Therefore the numerical rows are
conditional benchmarks and are not asserted as certified resistance
of the experimental generator.

\paragraph{Why the range cannot be omitted.}
There is an explicit contrasting calculation within the same circuit
family. Suppose the Euler angles are instead independent uniform
variables on $[0,\pi)$, independently of $B$. Averaging the Bloch
rotations gives, for a data qubit initially in $|0\rangle$ or
$|1\rangle$,
\begin{equation}
 \langle X\rangle_0=(2/\pi)^2,\qquad
 \langle X\rangle_1=-(2/\pi)^2.
\end{equation}
Consider all-black and all-white QIRBP images. At output position
$t=0$, no position control is active, so $V_B(0)=I$. That position
occurs with probability $1/N$, regardless of the permutation.
An adversary can request fresh encryptions until measuring $t=0$,
then measure $X$ on one colour-data qubit. This procedure uses $N$
preparations on average and succeeds with probability
\begin{equation}
 \frac{1+(2/\pi)^2}{2}\simeq0.702642>2/3.
\end{equation}
It is a conditional countermodel, not a claimed attack on an
unspecified nD-QHCS law. It shows mathematically why one cannot
assign a source-independent numerical resistance to the three-layer
architecture or to the statement that keys are refreshed.

\paragraph{Interpretation of the original attack labels.}
Independent known/chosen-image pairs from other operations do not
share the target key under the independent-refresh assumption. They
can still reveal a biased generator law or shared source parameters;
the conditional key law must therefore be specified. Generic key
search, if a reversible candidate verifier is available, has the
usual conditional $\Theta(\sqrt{M/t})$ query baseline for $M$ candidates
and $t$ marked candidates \cite{GDO,Boyer1998QEF,Zalka1999QEF}.
No particular old key-length value is used to certify QEF, and no
reduction proving that all attacks require exhaustive search is
claimed. The dimension or parameter count of a general quantum state
or unitary is not substituted for an attack lower bound.

Theorem~\ref{thm:qef_moment_sample_bound} concerns the mathematical
quantum operators and the stated key law. Classical array simulations
and their rendered images do not verify these quantum moment
conditions. In particular, the original scalar-seeded XOR/permutation
simulation is not an implementation of the controlled complex
Euler rotations assumed here. Image-quality and decryption metrics
are therefore reported as diagnostics, rather than as evidence that
the conditional security bound holds. The new probabilistic bound
is proved analytically; it is not included in the existing Lean~4
unitarity checks. Entanglement alone supplies neither authentication
nor a proof of resistance to active modification.
\par}
